% (User input window)
% 2026-02-15 21:36:22（Asia/Singapore）

\paragraph{纤维多重度的相对熵坐标：算术--几何缺口、端点压力与相变指数}
以下条目以固定分辨率 \(m\) 的纤维多重度账本为对象，将“边界体积 \(\log|X_m|\)”与“纤维对数体积 \(\log d_m\)”之间的差异精确压缩为 KL/Jeffreys 散度，
并由此给出三种平均的序链、两种全息分解、端点压力斜率以及线性增长率的临界指数化表达。

\begin{conclusion}[算术体积—几何体积的 KL 精确差：边界维度的熵型缺口恒等式]\label{con:xi-kl-arith-geo-gap-identity}
在分辨率 \(m\) 上令 \(\Fold_m:\Omega_m\to X_m\) 的纤维多重度为 \(d_m(x):=\card{\Fold_m^{-1}(x)}\)，并满足守恒
\(
\sum_{x\in X_m} d_m(x)=\card{\Omega_m}=2^m
\)。
在 \(X_m\) 上定义两分布
\[
u_m(x):=\frac{1}{|X_m|},
\qquad
w_m(x):=\frac{d_m(x)}{2^m}.
\]
则相对熵满足严格恒等式
\[
\boxed{\;
D_{\mathrm{KL}}(u_m\mid w_m)
=\log\frac{2^m}{|X_m|}
-\frac{1}{|X_m|}\sum_{x\in X_m}\log d_m(x)\;}
\]
从而
\(
\log(2^m/|X_m|)
\)
与
\(
|X_m|^{-1}\sum_x \log d_m(x)
\)
之差精确等于 \(u_m\) 到 \(w_m\) 的相对熵。
并且
\[
D_{\mathrm{KL}}(u_m\mid w_m)=0
\iff
u_m\equiv w_m
\iff
d_m(x)\equiv \frac{2^m}{|X_m|}\ \ (\forall x\in X_m).
\]
\end{conclusion}

\begin{conclusion}[维度—体积—熵的三层夹逼：\texorpdfstring{$\log d$}{log d} 的三种平均形成序链]\label{con:xi-logd-three-average-chain-kl}
沿用结论 \ref{con:xi-kl-arith-geo-gap-identity} 的记号，定义三种尺度指数
\[
\gamma_m^{\mathrm{geo}}
:=\frac{1}{m|X_m|}\sum_{x\in X_m}\log d_m(x),
\qquad
\gamma_m^{\mathrm{arith}}
:=\frac{1}{m}\log\frac{2^m}{|X_m|},
\qquad
\gamma_m^{\mathrm{Sh}}
:=\frac{1}{m2^m}\sum_{x\in X_m}d_m(x)\log d_m(x).
\]
则成立序链
\[
\gamma_m^{\mathrm{geo}}\ \le\ \gamma_m^{\mathrm{arith}}\ \le\ \gamma_m^{\mathrm{Sh}},
\]
并且两端差值具有精确的 KL 表达
\[
\gamma_m^{\mathrm{arith}}-\gamma_m^{\mathrm{geo}}=\frac{1}{m}D_{\mathrm{KL}}(u_m\mid w_m),
\qquad
\gamma_m^{\mathrm{Sh}}-\gamma_m^{\mathrm{arith}}=\frac{1}{m}D_{\mathrm{KL}}(w_m\mid u_m).
\]
因此两条不等式同时取等当且仅当纤维完全等多重（即 \(u_m\equiv w_m\)）。
\end{conclusion}

\begin{conclusion}[Jeffreys 散度即分形不对称的数值证书：两种可见平均之差的精确公式]\label{con:xi-jeffreys-logd-gap-certificate}
在结论 \ref{con:xi-logd-three-average-chain-kl} 的记号下，定义 Jeffreys 散度
\[
J(u_m,w_m):=D_{\mathrm{KL}}(u_m\mid w_m)+D_{\mathrm{KL}}(w_m\mid u_m).
\]
则有严格恒等式
\[
\gamma_m^{\mathrm{Sh}}-\gamma_m^{\mathrm{geo}}=\frac{1}{m}J(u_m,w_m).
\]
故“按边界均匀抽样”与“按微观均匀抽样（推前权重 \(d_m\)）”两种口径下看到的平均 \(\log d_m\) 之差，恰为 Jeffreys 散度的尺度化；
该量为零当且仅当 \(u_m\equiv w_m\)。
\end{conclusion}

\begin{conclusion}[微观熵密度的两种全息分解：体积项、边界项与相对熵修正]\label{con:xi-micro-entropy-holographic-two-splits}
仍令 \(\sum_{x\in X_m} d_m(x)=2^m\)，则二进制微观熵密度 \(m\log 2\) 具有两条互补的严格分解：
\[
m\log 2
=\log|X_m|
+\sum_{x\in X_m}u_m(x)\log d_m(x)\ +\ D_{\mathrm{KL}}(u_m\mid w_m),
\]
以及
\[
m\log 2
=\log|X_m|
+\sum_{x\in X_m}w_m(x)\log d_m(x)\ -\ D_{\mathrm{KL}}(w_m\mid u_m).
\]
因此同一“总熵”可被写成：边界体积项 \(\log|X_m|\) 与纤维对数体积的期望之和，再加上或减去一个相对熵修正项；
相对熵项精确度量了在非匹配测度下读取纤维对数体积时引入的系统偏置。
\end{conclusion}

\begin{conclusion}[等多重刚性判据：尺度化相对熵消失当且仅当几何均值指数贴住算术指数]\label{con:xi-equimultiplicity-kl-vanishing}
设极限
\[
h_\partial:=\lim_{m\to\infty}\frac{1}{m}\log|X_m|
\]
存在，且 \(\lim_{m\to\infty}\gamma_m^{\mathrm{geo}}\) 存在。
则以下条件等价：
\begin{enumerate}
  \item \(\displaystyle \lim_{m\to\infty}\frac{1}{m}D_{\mathrm{KL}}(u_m\mid w_m)=0\)；
  \item \(\displaystyle \lim_{m\to\infty}\gamma_m^{\mathrm{geo}}=\log 2-h_\partial\)；
  \item \(\displaystyle \lim_{m\to\infty}\bigl(\gamma_m^{\mathrm{arith}}-\gamma_m^{\mathrm{geo}}\bigr)=0\)。
\end{enumerate}
故尺度化相对熵是否消失可作为“边界厚度均匀性”的渐近判据：若其消失，则几何均值指数达到由边界体积项强制的上界；若不消失，则存在不可压缩的熵型缺口。
\end{conclusion}

\begin{conclusion}[对数体积的两种极限端点：\texorpdfstring{$q\to 0$}{q->0} 与 \texorpdfstring{$q\to 1$}{q->1} 的端点导数]\label{con:xi-pressure-endpoints-geo-shannon}
令
\(
S_q(m):=\sum_{x\in X_m} d_m(x)^q
\)，并设对每个固定 \(q>0\) 极限
\[
\tau(q):=\lim_{m\to\infty}\frac{1}{m}\log S_q(m)
\]
存在。
则在存在一侧导数的意义下，端点导数满足
\[
\tau'(0^+)=\lim_{m\to\infty}\gamma_m^{\mathrm{geo}},
\qquad
\tau'(1^-)=\lim_{m\to\infty}\gamma_m^{\mathrm{Sh}}.
\]
因此同一压力函数 \(\tau\) 的两端斜率分别编码“边界均匀口径”的几何均值体积指数与“微观均匀口径”的 Shannon 加权体积指数。
\end{conclusion}

\begin{conclusion}[中心规范熵的普适剥离：\texorpdfstring{$(\ZZ_2)^b$}{(Z2)^b} 的对角约束导致体积亏损]\label{con:xi-gauge-center-diagonal-strip}
设存在 \(b=b(m)\) 条纤维满足 \(d_m(x)=2\)，并令这些 \(S_2\) 因子在纤维对称群直积中生成中心子群 \((\ZZ_2)^b\)。
若进一步施加全局几何可实现约束，使得仅允许对角翻转子群 \(\ZZ_2^{\mathrm{diag}}\subset(\ZZ_2)^b\) 保留，则允许的中心体积发生严格剥离
\[
|(\ZZ_2)^b|/|\ZZ_2^{\mathrm{diag}}|=2^{b-1},
\qquad
\Delta S_{\mathrm{gauge}}=(b-1)\log 2.
\]
当 \(b-1\) 为偶数时，上述体积因子可改写为 \(4^{(b-1)/2}\)，从而在对数账本中呈现强制的 \(4\)-进量子化层级。
\end{conclusion}

\begin{conclusion}[边界相变的熵学判据：两种 KL 的线性增长率给出临界指数]\label{con:xi-boundary-phase-transition-kl-rates}
若存在常数 \(c_{u\to w},c_{w\to u}\ge 0\) 使
\[
\lim_{m\to\infty}\frac{1}{m}D_{\mathrm{KL}}(u_m\mid w_m)=c_{u\to w},
\qquad
\lim_{m\to\infty}\frac{1}{m}D_{\mathrm{KL}}(w_m\mid u_m)=c_{w\to u},
\]
并且 \(h_\partial=\lim_{m\to\infty}m^{-1}\log|X_m|\) 存在，
则由结论 \ref{con:xi-logd-three-average-chain-kl}--\ref{con:xi-jeffreys-logd-gap-certificate} 得到
\[
\lim_{m\to\infty}\gamma_m^{\mathrm{geo}}=\log 2-h_\partial-c_{u\to w},
\qquad
\lim_{m\to\infty}\gamma_m^{\mathrm{Sh}}=\log 2-h_\partial+c_{w\to u},
\]
以及
\[
\lim_{m\to\infty}\bigl(\gamma_m^{\mathrm{Sh}}-\gamma_m^{\mathrm{geo}}\bigr)=c_{u\to w}+c_{w\to u}.
\]
因此可将 \(c_{u\to w}+c_{w\to u}\) 视为“边界相变强度”的一元证书：它同时度量两端平均的分离幅度，并等于 Jeffreys 散度的线性增长率。
\end{conclusion}

\begin{conclusion}[分形维度的熵化下界：维度密度被 \texorpdfstring{$c_{u\to w}$}{c_u->w} 严格压低]\label{con:xi-fractal-dimension-entropy-deficit}
设 \(\varphi\)-标度下的边界维度密度
\[
d_\partial:=\lim_{m\to\infty}\frac{1}{m\log\varphi}\log|X_m|
\]
存在。
在结论 \ref{con:xi-boundary-phase-transition-kl-rates} 的假设下，纤维几何均值的 \(\varphi\)-维度密度
\[
d_{\mathrm{fib}}^{\mathrm{geo}}
:=\frac{1}{\log\varphi}\lim_{m\to\infty}\gamma_m^{\mathrm{geo}}
\]
满足严格表达
\[
d_{\mathrm{fib}}^{\mathrm{geo}}=\frac{\log 2}{\log\varphi}-d_\partial-\frac{c_{u\to w}}{\log\varphi}.
\]
因此一旦 \(c_{u\to w}>0\)，则 \(d_\partial+d_{\mathrm{fib}}^{\mathrm{geo}}\) 将严格低于微观上限 \(\log 2/\log\varphi\)，其严格亏损由相对熵增长率封口。
\end{conclusion}

\begin{conclusion}[熵—体积的可逆重建：\texorpdfstring{$(D_{\mathrm{KL}},\gamma)$}{(DKL,gamma)} 的双坐标互定]\label{con:xi-entropy-volume-dual-coordinates}
在固定 \(m\) 上给定 \(\log|X_m|\) 与总量 \(2^m\)。
则两组数据互相可逆确定：
\begin{itemize}
  \item 熵坐标：\(\bigl(D_{\mathrm{KL}}(u_m\mid w_m),\,D_{\mathrm{KL}}(w_m\mid u_m)\bigr)\)；
  \item 体积坐标：\(\bigl(\gamma_m^{\mathrm{geo}},\,\gamma_m^{\mathrm{Sh}}\bigr)\)。
\end{itemize}
具体地，记
\[
\gamma_m^{\mathrm{arith}}
=\frac{1}{m}\log\frac{2^m}{|X_m|}
=\log 2-\frac{1}{m}\log|X_m|
\]
为已知，则由结论 \ref{con:xi-logd-three-average-chain-kl} 得到双向线性重建
\[
D_{\mathrm{KL}}(u_m\mid w_m)=m\bigl(\gamma_m^{\mathrm{arith}}-\gamma_m^{\mathrm{geo}}\bigr),
\qquad
D_{\mathrm{KL}}(w_m\mid u_m)=m\bigl(\gamma_m^{\mathrm{Sh}}-\gamma_m^{\mathrm{arith}}\bigr),
\]
等价地
\[
\gamma_m^{\mathrm{geo}}=\gamma_m^{\mathrm{arith}}-\frac{1}{m}D_{\mathrm{KL}}(u_m\mid w_m),
\qquad
\gamma_m^{\mathrm{Sh}}=\gamma_m^{\mathrm{arith}}+\frac{1}{m}D_{\mathrm{KL}}(w_m\mid u_m).
\]
因此在该纤维账本框架内，“熵”并非附加量，而是对数体积的等价坐标系；边界厚度与端点可见平均的全部一阶信息可在该坐标变换下等价转写。
\end{conclusion}

